\section*{Bab \ref{ch:g12:selection-drift-speciation} --- Seleksi, Hanyutan dan Spesiasi}

\begin{solution}{exo:g12:selection-drift-speciation:1}
Bagian dari semua salinan sebuah \omterm{def:g10:universal-dna:gene}{gen} dalam \omterm{def:g12:selection-drift-speciation:population}{populasi} yang berupa \omterm{def:g10:universal-dna:gene}{alel}
tertentu. Dengan $p = 0.8$, $q = 0.2$: $AA$ 64\%, $Aa$ 32\%, $aa$ 4\%.
\end{solution}

\begin{solution}{exo:g12:selection-drift-speciation:2}
Seleksi: perbedaan reproduksi antara \omterm{def:g11:genes-and-disease:genetic}{pembawa} \omterm{def:g10:universal-dna:gene}{alel} yang berbeda, yang
ditetapkan oleh lingkungan, dan mengubah frekuensi ke satu arah.
Hanyutan: ayunan acak frekuensi akibat penyampelan \omterm{def:g12:selection-drift-speciation:population}{populasi} terhingga,
tanpa arah. Seleksi bergantung pada apa yang dikerjakan \omterm{def:g10:universal-dna:gene}{alel} itu;
hanyutan tidak.
\end{solution}

\begin{solution}{exo:g12:selection-drift-speciation:3}
Sekitar 75\% pada 1880, ketika jelaga telah menghitamkan batang pohon;
sekitar 30\% pada 1980, setelah undang-undang udara bersih membiarkan
batang itu memucat lagi.
\end{solution}

\begin{solution}{exo:g12:selection-drift-speciation:4}
\omterm{def:g10:universal-dna:gene}{Alel} generasi berikutnya adalah sampel dari generasi sebelumnya;
sampel kecil menyimpang lebih jauh dari sumbernya daripada sampel
besar, sehingga frekuensinya berayun lebih lebar, dan sebuah \omterm{def:g10:universal-dna:gene}{alel}
dapat hilang oleh kebetulan.
\end{solution}

\begin{solution}{exo:g12:selection-drift-speciation:5}
Terputusnya pertukaran \omterm{def:g10:universal-dna:gene}{gen} antara dua \omterm{def:g12:selection-drift-speciation:population}{populasi}, cukup lama sehingga
keduanya memencar sampai tak dapat lagi saling kawin.
\end{solution}

\begin{solution}{exo:g12:selection-drift-speciation:6}
Salinan $a$: $160 + 40 = 200$ dari 1000, sehingga $q = 0.2$, $p = 0.8$.
Harapan: $0.64 \times 500 = 320$, $0.32 \times 500 = 160$, $0.04
\times 500 = 20$ --- persis sesuai sensus: \omterm{def:g12:selection-drift-speciation:population}{populasi} itu berada pada
keadaan acuan.
\end{solution}

\begin{solution}{exo:g12:selection-drift-speciation:7}
$q^2 = 0.0001$: satu dari sepuluh ribu menampakkan sifat itu; $2pq \approx
0.02$: satu dari lima puluh membawanya. Sembilan puluh sembilan persen
salinannya duduk di dalam heterozigot, yang tak tersentuh oleh seleksi
terhadap sifat itu.
\end{solution}

\begin{solution}{exo:g12:selection-drift-speciation:8}
Keduanya mencapai ujung pada generasi 12, satu terpancang di 100\%,
satunya hilang. Yang ketiga mengembara tanpa menyentuh batas mana pun
dalam 20 generasi --- kebetulan lagi; diberi waktu lebih, ia akan
sampai juga.
\end{solution}

\begin{solution}{exo:g12:selection-drift-speciation:9}
Ngengat pucat terus berbiak di pedesaan yang tak tercemar, tempat
mereka diuntungkan, dan ngengat itu terbang: pendatang membawa \omterm{def:g10:universal-dna:gene}{alel}
pucat kembali ke wilayah berjelaga setiap generasi, sehingga \omterm{def:g10:universal-dna:gene}{alel} itu
tak pernah lenyap.
\end{solution}

\begin{solution}{exo:g12:selection-drift-speciation:10}
Dua puluh individu membawa paling banyak empat puluh salinan tiap \omterm{def:g10:universal-dna:gene}{gen},
sampel yang mungil dari \omterm{def:g10:universal-dna:gene}{alel} aslinya; sebagian besar hilang dalam
leher botol itu, dan hanyutan dalam \omterm{def:g12:selection-drift-speciation:population}{populasi} kecil sesudahnya
menghilangkan lebih banyak lagi. Jumlahnya kembali; \omterm{def:g10:universal-dna:gene}{alelnya} tidak,
kecuali oleh \omterm{def:g11:mutations:mutation}{mutasi} baru.
\end{solution}

\begin{solution}{exo:g12:selection-drift-speciation:11}
Sepuluh \omterm{def:g12:selection-drift-speciation:population}{populasi} yang bergerak searah, setelah perubahan lingkungan
yang sama, dengan mekanisme yang masuk akal: seleksi. Satu \omterm{def:g12:selection-drift-speciation:population}{populasi}
mungil berisi tiga puluh ekor, tanpa sebab yang dilaporkan: hanyutan
adalah penjelasan bakunya sampai sebuah mekanisme ditunjukkan.
\end{solution}

\begin{solution}{exo:g12:selection-drift-speciation:12}
$AA$ $0.85^2 \approx 72\%$, $Aa$ $2 \times 0.85 \times 0.15 \approx
26\%$, $aa$ $0.15^2 \approx 2\%$. Yang $aa$ mati, menyingkirkan salinan
$a$; tetapi yang $Aa$, seperempat \omterm{def:g12:selection-drift-speciation:population}{populasi}, bertahan terhadap malaria
lebih baik daripada $AA$ dan mewariskan salinan $a$ mereka. Kehilangan
dari $aa$ diimbangi oleh keuntungan $Aa$, dan $a$ bertahan di 15\%.
\end{solution}

\begin{solution}{exo:g12:selection-drift-speciation:13}
Pemilihan pasangan berdasarkan warna: sebuah penghalang perilaku.
Dalam air keruh penghalang itu gagal dan kedua \omterm{def:g10:biodiversity-scales:species}{spesies} berhibrida,
sehingga belum ada ketakserasian genetik yang terhimpun: pemisahannya
baru dan masih dapat berbalik.
\end{solution}

\begin{solution}{exo:g12:selection-drift-speciation:14}
Sepanjang cincin itu setiap \omterm{def:g12:selection-drift-speciation:population}{populasi} saling kawin dengan tetangganya,
dengan perbedaan kecil saja antara yang bersebelahan; perbedaan itu
menumpuk mengelilingi lingkaran sampai kedua ujungnya, ketika bertemu,
tidak saling kawin. Seluruh gradiennya terlihat sekaligus: spesiasi
adalah soal langkah-langkah kecil yang terhimpun. Kedua ujung itu
disebut \omterm{def:g10:biodiversity-scales:species}{spesies} karena, di tempat mereka bertemu, keduanya berlaku
sebagai \omterm{def:g10:biodiversity-scales:species}{spesies}: tanpa pertukaran \omterm{def:g10:universal-dna:gene}{gen}.
\end{solution}

\begin{solution}{exo:g12:selection-drift-speciation:15}
Seleksi membuat \omterm{def:g11:genes-and-disease:genetic}{pembawa} sebagian \omterm{def:g10:universal-dna:gene}{alel} bereproduksi lebih banyak
\emph{di sini dan sekarang}: lebih baik dalam meninggalkan keturunan
di lingkungan ini, dibandingkan dengan anggota lain \omterm{def:g12:selection-drift-speciation:population}{populasinya} ---
bukan lebih baik dalam arti mutlak apa pun. Ketika lingkungan berubah,
arahnya berbalik (ngengat itu); hanyutan membuat \omterm{def:g12:selection-drift-speciation:population}{populasi} berbeda
tanpa alasan adaptasi; seleksi seksual menyukai perhiasan yang justru
menghambat kelangsungan hidup. Evolusi tidak punya arah perbaikan,
hanya pemilahan setempat.
\end{solution}

\begin{solution}{pb:g12:selection-drift-speciation:1}
\textbf{1.} Salinan $D$: $2 \times 720 + 960 = 2400$; salinan $d$: $960 +
2 \times 320 = 1600$; dari 4000: $p = 0.6$, $q = 0.4$.

\textbf{2.} Harapan $0.36 \times 2000 = 720$, $0.48 \times 2000 =
960$, $0.16 \times 2000 = 320$: persis sesuai sensus. Ya.

\textbf{3.} $960/1680 \approx 57\%$.

\textbf{4.} Sebagian besar salinan $d$ ada pada \omterm{def:g11:genes-and-disease:genetic}{pembawa} gelap. Setelah
penyingkiran: 1680 tikus, salinan $D$ 2400, salinan $d$ 960: $q = 960/3360 \approx
0.29$ --- turun dari 0.40, tetapi jauh dari lenyap.

\textbf{5.} Ukuran yang besar: \omterm{def:g12:selection-drift-speciation:population}{populasi} pulau yang kecil menghanyut,
sehingga frekuensinya mengembara oleh kebetulan (migrasi dari daratan
utama tidak ada, \omterm{def:g11:mutations:mutation}{mutasi} dapat diabaikan).

\textbf{6.} $DD$ 360, $Dd$ 480, $dd$ 320: 1160 ekor pembiak.

\textbf{7.} $D$: $720 + 480 = 1200$; $d$: $480 + 640 = 1120$; dari 2320:
$p \approx 0.52$, $q \approx 0.48$.

\textbf{8.} Generasi berikutnya (per 2000): $DD$ $0.52^2 \approx 27\%$
(535), $Dd$ $\approx 50\%$ (999), $dd$ $\approx 23\%$ (466). Setelah
burung hantu: 267, 500, 466. $D$: $534 + 500 = 1034$; $d$: $500 + 932 =
1432$; $q \approx 0.58$.

\textbf{9.} Dengan $q = 0.58$: $DD$ 18\% (353), $Dd$ 49\% (974), $dd$
34\% (673); penyintas 176, 487, 673; $d$: $487 + 1346 = 1833$ dari 2672:
$q \approx 0.69$. Kecenderungannya: 0.40, 0.48, 0.58, 0.69 --- naik sekitar
0.1 tiap generasi.

\textbf{10.} Setiap \omterm{def:g11:genes-and-disease:genetic}{pembawa} $D$ berbulu gelap dan terpapar burung
hantu, sehingga seleksi bekerja pada setiap salinan $D$ pada setiap
generasi; $q$ naik terus. \omterm{def:g10:universal-dna:gene}{Alel} \omterm{def:g11:genes-and-disease:genetic}{resesif} akan bersembunyi di dalam
heterozigot begitu menjadi jarang, tetapi \omterm{def:g10:universal-dna:gene}{alel} \omterm{def:g11:genes-and-disease:genetic}{dominan} tak pernah
bersembunyi: $D$ dapat didorong sampai nol.

\textbf{11.} $D$: $4 + 3 = 7$; $d$: $3 + 2 = 5$; $q = 5/12 \approx
0.42$, dekat dengan 0.40 pulau besar oleh untung-untungan.

\textbf{12.} Di antara para pembiak, $D$ punya 5 salinan dan $d$ 1: $q$ pada
generasi berikutnya rata-rata $1/6$, tetapi salinan tunggal $d$ itu
mungkin tidak diwariskan kepada satu anak pun ($q = 0$) atau kepada
beberapa; \omterm{def:g11:enzymes-and-phenotype:phenotype}{genotipe} $dd$ sudah lenyap, dan bersamanya sebagian besar
keragamannya.

\textbf{13.} Dengan pembiak sesedikit itu, \omterm{def:g10:universal-dna:gene}{alel} mana yang diwariskan
adalah soal kebetulan pada setiap generasi; \omterm{def:g10:universal-dna:gene}{alel} yang jarang dapat
hilang dalam satu anakan yang bernasib buruk.

\textbf{14.} $d$ akan menghanyut, dan dalam beberapa generasi menjadi
terpancang atau hilang --- kedua hasil itu oleh kebetulan. Pulau kecil
tetangga, dari titik awal yang sama, dapat berakhir dengan \omterm{def:g10:universal-dna:gene}{alel} yang
berlawanan: kedua pulau berbeda warna bulu tanpa alasan adaptif.

\textbf{15.} \omterm{prop:g12:selection-drift-speciation:drift}{Hanyutan genetik}; ukuran koloni yang mungil.

\textbf{16.} Ya: mereka jarang kawin dengan tikus pulau besar dan
hibridanya mandul --- isolasi reproduktif.

\textbf{17.} Lautnya, sebuah penghalang geografi; seleksi di
lingkungan yang berbeda (tanpa burung hantu, dengan kondisi lain) dan
\omterm{prop:g12:selection-drift-speciation:drift}{hanyutan genetik} dalam \omterm{def:g12:selection-drift-speciation:population}{populasi} kecil.

\textbf{18.} Tikus yang berbiak pada musim berbeda tak pernah bertemu
sebagai pasangan: bahkan tanpa penghalang geografi apa pun, tak ada
\omterm{def:g10:universal-dna:gene}{gen} yang mengalir, dan kedua \omterm{def:g12:selection-drift-speciation:population}{populasi} terus memencar.

\textbf{19.} Tidak sekarang: keduanya terisolasi oleh perilaku dan
kemandulan, dan akan tetap menjadi dua \omterm{def:g10:biodiversity-scales:species}{spesies} berdampingan. Setelah
100 tahun saja keduanya masih akan saling kawin dengan bebas dan
melebur kembali menjadi satu \omterm{def:g12:selection-drift-speciation:population}{populasi}.

\textbf{20.} $q$ dari 0.40 menjadi sekitar 0.69 setelah tiga generasi
burung hantu; di pulau kecil, $d$ terpancang atau hilang oleh \omterm{prop:g12:selection-drift-speciation:drift}{hanyutan
genetik}; dua \omterm{def:g10:biodiversity-scales:species}{spesies} pada akhirnya.
\end{solution}
